\section{机器人移动：腿式运动}

\subsection{四足机器人行走}

讲者展示了使用 RL 训练四足机器人（如 Unitree Go1/A1）在各种地形上行走的案例。

\begin{importantbox}{腿式运动的 RL 框架}
\begin{itemize}
    \item \textbf{状态}：关节角度、角速度、IMU 数据（倾斜角、角速度）
    \item \textbf{动作}：每个关节的目标角度（PD 控制器跟踪）
    \item \textbf{奖励}：前进速度 + 能量惩罚 + 姿态惩罚 + 脚步频率奖励
    \item \textbf{训练}：在 Isaac Gym 等 GPU 并行仿真器中进行，可以同时运行数千个环境
\end{itemize}
\end{importantbox}

\subsection{Teacher-Student 框架}

\begin{knowledgebox}{特权信息蒸馏}
一种强大的 sim-to-real 训练范式：
\begin{enumerate}
    \item \textbf{Teacher 阶段}：在仿真中训练 teacher 策略，允许访问\textbf{特权信息}（如精确的地形高度图、物体质量、摩擦系数）
    \item \textbf{Student 阶段}：训练 student 策略仅使用\textbf{真实可用的传感器}（如本体感知、摄像头），通过蒸馏模仿 teacher 的行为
\end{enumerate}
Teacher 可以学得更好（因为有更多信息），student 则学会从有限传感器中推断出等效的信息。
\end{knowledgebox}

\begin{figure}[H]
\centering
\includegraphics[width=0.9\textwidth]{slides-images/slide-16.jpg}
\caption{部署阶段的关键不是“直接复用仿真策略”，而是根据历史观测在线估计外参并持续校正控制器，这也是现代腿式机器人鲁棒性的核心来源。}
\end{figure}

\subsection{本章小结}

RL + sim-to-real 已经成为腿式机器人控制的主流范式，teacher-student 框架有效解决了感知受限问题。

